\documentclass[11pt]{article}
\input{style-pc}
\usepackage{array}

\begin{document}

\entetesujet{Sujet bac 2011 -- Physique-Chimie -- Série D}

% ============================ CHIMIE ============================
\matiere{CHIMIE}{8 points}

\exo{1}{}
La réaction de décomposition de $\ch{NOBr}$, à une température déterminée, selon l'équation $\ch{NOBr_{(g)} \longrightarrow 2\,NO_{(g)} + Br_{2(g)}}$ fournit les résultats suivants :

\begin{center}\renewcommand{\arraystretch}{1.3}
\begin{tabular}{|l|c|c|c|c|c|c|}
\hline
$t$(s) & 0 & 6,2 & 10,8 & 14,7 & 20,0 & 24,5 \\
\hline
$[\ch{NOBr}]$ mol$\cdot$L$^{-1}$ & 0,0250 & 0,0198 & 0,0162 & 0,0144 & 0,0125 & 0,0112 \\
\hline
\end{tabular}
\end{center}

\begin{enumerate}
  \item En se servant de ces résultats, déterminer le temps de demi-réaction.
  \item Sachant que la réaction est d'ordre deux :
  \begin{enumerate}[label=\alph*.]
    \item Calculer la constante de vitesse de la réaction.
    \item Écrire la loi de vitesse de cette réaction.
    \item Déterminer le temps nécessaire à la disparition de 80 \% du réactif, puis calculer la vitesse de disparition du réactif à cette date.
  \end{enumerate}
\end{enumerate}

\exo{2}{}
On dissout $2{,}3$ g d'acide méthanoïque $\ch{HCOOH}$ dans l'eau pure de façon à obtenir 500 mL de solution. Toutes les mesures sont réalisées à $25\,^\circ$C. Le pH de cette solution est $2{,}4$.

\begin{enumerate}
  \item \begin{enumerate}[label=\alph*.]
    \item Calculer la concentration molaire volumique de la solution préparée.
    \item Montrer que l'acide méthanoïque est un acide faible.
    \item Écrire l'équation de dissociation de l'acide méthanoïque dans l'eau.
  \end{enumerate}
  \item \begin{enumerate}[label=\alph*.]
    \item Faire l'inventaire de toutes les espèces chimiques présentes dans la solution.
    \item Calculer les concentrations molaires volumiques des différentes espèces chimiques.
  \end{enumerate}
  \item \begin{enumerate}[label=\alph*.]
    \item Calculer le pKa du couple $\ch{HCOOH/HCOO^-}$.
    \item Le pKa du couple $\ch{CH_3COOH/CH_3COO^-}$ est $4{,}7$. Comparer les forces des acides méthanoïque ($\ch{HCOOH}$) et éthanoïque ($\ch{CH_3COOH}$).
  \end{enumerate}
\end{enumerate}
\noindent On donne en g/mol : C : 12 ; O : 16 ; H : 1.

% ============================ PHYSIQUE ============================
\matiere{PHYSIQUE}{12 points}

\exo{1}{}
On considère un ressort élastique de masse négligeable, de constante de raideur $K = 26$ N$\cdot$m$^{-1}$, suspendu verticalement par l'une des extrémités à une potence. À l'autre extrémité, on fixe un solide $(S)$ de masse $m$ ; le ressort s'allonge de $\Delta l_0$.

\begin{enumerate}
  \item Écrire la relation donnant l'allongement $\Delta l_0$ du ressort à l'équilibre en fonction de $k$, $m$ et $g$.
  \item Le solide $(S)$ est écarté de sa position d'équilibre de 3 cm vers le bas, puis lâché sans vitesse initiale à la date $t = 0$. La période des oscillations libres est $T = 0{,}52$ s.
  \begin{enumerate}[label=\alph*.]
    \item Établir l'équation différentielle du mouvement du solide $(S)$.
    \item Déterminer la masse de ce solide.
  \end{enumerate}
  \item Déterminer l'équation horaire du mouvement de $(S)$.
  \item Trouver la vitesse du solide $(S)$ au premier passage par la position d'équilibre.
\end{enumerate}
\noindent On donne $g = 10$ m/s$^2$.

\exo{2}{}
L'extrémité $A$ d'une corde élastique est animée d'un mouvement vibratoire dont l'élongation instantanée exprimée en mètres est $Y_A = 4\cdot10^{-2}\sin 20\pi t$.

\begin{enumerate}
  \item Déterminer l'amplitude, la période, la fréquence et la phase initiale du mouvement de $A$.
  \item La célérité du mouvement vibratoire est $c = 2{,}5$ m/s. Déterminer :
  \begin{enumerate}[label=\alph*.]
    \item La longueur d'onde du mouvement vibratoire.
    \item L'équation horaire du mouvement d'un point $M$ de la corde situé à une distance $d = 62{,}5$ cm de $A$.
    \item Représenter le graphe du mouvement de $M$.
  \end{enumerate}
  \item On considère un point $N$ situé à $93{,}75$ cm de $A$. Comparer les mouvements de $M$ et $N$ à celui de $A$.
\end{enumerate}

\exo{3}{}
\begin{enumerate}
  \item Un conducteur ohmique de résistance $R = 20$ $\Omega$ est parcouru par un courant alternatif sinusoïdal de fréquence $N = 50$ Hz, d'intensité instantanée $i(t) = 2\sqrt{2}\cos\omega t$. Donner l'expression de la tension instantanée aux bornes du conducteur ohmique.
  \item On monte en série avec le conducteur ohmique précédent un condensateur de capacité $C = 2\cdot10^{-4}$ F. L'ensemble est parcouru par le courant alternatif précédent.
  \begin{enumerate}[label=\alph*.]
    \item Faire le schéma du circuit.
    \item Calculer l'impédance du circuit ainsi constitué.
    \item Déterminer le déphasage de la tension aux bornes des deux dipôles par rapport à l'intensité.
    \item Établir l'expression de la tension instantanée $u(t)$.
  \end{enumerate}
  \item Calculer la puissance moyenne consommée dans le circuit.
\end{enumerate}

\end{document}
